\documentclass[english, parskip=full]{scrartcl}
\usepackage[utf8]{inputenc}  % Kodierung der Datei
\usepackage[english]{babel}  % Sprache auf Deutsch 
\usepackage{color}
\usepackage{graphicx, gensymb}
\usepackage{amsfonts, cancel, amssymb}
\usepackage{amsmath, amsthm, array, esint}
\usepackage{booktabs, multirow}  % schönere Tabellen
\usepackage{caption}  % Anpassung der Captions
\usepackage{scrlayer-scrpage}    % Manipulation des Seitenstils
\pagestyle{scrheadings}  % Stil für die Seite setzen
\automark{section}  % Abschnittsnamen als Seitenbeschriftung 
\setheadsepline{.5pt}    % Kopzeile durch Linie abtrennen
\usepackage[hidelinks]{hyperref}  % Links und weitere PDF-Features
\usepackage{typearea, tikz, pgffor, verbatim}
\usetikzlibrary{calc, arrows.meta,arrows, graphs, quotes, positioning}
\usepackage[cache=false, outputdir=build]{minted}
\usemintedstyle{vs}
\usepackage{placeins} % provides \FloatBarrier

\definecolor{bg}{rgb}{0.9,0.9,0.9}
\newcommand\numberthis{\addtocounter{equation}{1}\tag{\theequation}}
\usepackage{svg}
\usepackage{siunitx}
\usepackage{physics}
\usepackage{tensor}
\renewcommand{\bra}{}
\newcommand{\kom}{}
\newcommand{\brak}{}
\renewcommand{\braket}{}
\newcommand{\intx}{}
\newcommand{\intr}{}
\newcommand{\kla}{}
\newcommand{\klam}{}
\newcommand{\klammer}{}
\newcommand{\oper}{}
\newcommand{\tecxt}{}
\newcommand{\Bra}{}
\newcommand{\Ket}{}
\renewcommand{\i}{\text{i}}
\newcommand{\e}{\text{e}}
\newcommand*{\Fourier}{\textsc{Fourier}\ }
\newcommand*{\F}{\mathcal{F}}
\newcommand*{\sinc}{\text{sinc\,}}
\newcommand{\conv}{\mathop{\scalebox{3.5}{\raisebox{-0.38ex}{\makebox[0.5\width][c]{$\ast$}}}}\limits}

\newcommand*{\titel}{Background and Central Theorems}
\newcommand*{\autor}{Joachim Schwardt}

\hypersetup{pdfauthor={\autor}, pdftitle={\titel}}

%\titlehead{titlehead}
\subject{Extreme Value Theory}
\title{\titel}
\author{\autor}
\date{\today}

\begin{document}

%\begin{titlepage}
\maketitle  % Titel setzen
%\tableofcontents  % Inhaltsverzeichnis setzen
%\end{titlepage}

\section{Asymptotic Extreme Value Distributions}
For a set of $n$ identically and independently distributed (iid) random variables $\xi_1,\dots,\xi_n$ consider the maximum value
\begin{align}
M_n &:= \max \klammer\left\{ \xi_j \right\}.
\end{align}
In general, the cumulative distribution is simply given by
\begin{align}
P(M_n\le x) &= P(\xi_1 \le x, \dots, \xi_n\le x) = F(x)^n
\end{align}
where $F(x)$ is the cumulative distribution function of the individual $\xi_j$.
The distribution of this maximum is governed only by the behaviour of $f(x) = F'(x)$ for large $x$, and locks the limiting distribution of $M_n$ into one of three possible forms.\\
For sequences of numbers $a_n>0, b_n$ consider the behaviour of $P(a_n(M_n - b_n) \le x)$ for large $n$.
Then the only possible limiting \emph{Extreme Value Distributions}, or \emph{max-stable} distributions, are
\begin{align}
\tecxt\text{Type I:}\quad G(x) &= \e^{-\e^{-x}}, &&& (\tecxt\text{Gumball class})\\
\tecxt\text{Type II:}\quad G(x) &= \Theta(x) \e^{-x^{-\alpha}},  &&& (\tecxt\text{Fr\'echet class})\\
\tecxt\text{Type III:}\quad G(x) &= \Theta(-x) \e^{-(-x)^\alpha} + \Theta(x), &&& (\tecxt\text{Weibull class})
\end{align}
where $\alpha > 0$.
These functions are also shown for some $\alpha$ in \autoref{fig:ev distributions simple}.
\begin{figure}[!ht]
\centering
\includegraphics{../EE_Pictures/ev_distributions_simple.png}
\caption{Plot of the three fundamental EV distributions for a fixed $\alpha$.}
\label{fig:ev distributions simple}
\end{figure}
\subsection{Khintchine's Convergence Theorem}
We consider the convergence of sequences
\begin{align}
F\kla\left( \frac{x}{a_n} + b_n \right) \overset{n\rightarrow \infty}{\longrightarrow} G(x)
\end{align}
and say that $F$ belongs to the \emph{domain of attraction} of $G$ (written $F\in D(G)$).
Before we continue with the convergence theory of these distributions, consider the inverses of continuous monotonically increasing functions given by
\begin{align}
\psi^{-1}(y) &= \inf \klammer\left\{ x \,\middle\vert\, \psi(x) \ge y \right\}.
\end{align}
Then for $a>0,b,c\in\mathbb{R}$ and $H(x) := \psi(ax+b) - c$ we find $H^{-1}(y) = \frac{1}{a} [\psi^{-1}(y+c) - b]$.
Khintchine's Theorem states that given a sequence of distribution functions $\{F_n\}$, $a_n>0,b_n$ and a non-degenerate d.f. $G$ for which
\begin{align}
F_n(a_nx+b_n)&\rightarrow G_1(x),
\end{align}
then for some $\alpha_n>0,\beta_n$ we find 
\begin{align}
F_n(\alpha_n x+ \beta_n) \rightarrow G_2(x)
\end{align}
if and only if $\frac{\alpha_n}{a_n} \rightarrow a$ and $\frac{\beta_n - b_n}{a_n} \rightarrow b$ for some $a>0,b$, and then $G_2(x) = G_1(ax+b)$.
\begin{proof}
Write $\tilde{F}_n(x) := F_n(a_nx + b_n)$, $\tilde{\alpha}_n := \frac{\alpha_n}{a_n}$ and $\tilde{\beta}_n := \frac{\beta_n - b_n}{a_n}$.
Then
\begin{align}
\tilde{F}_n(x) \rightarrow G_1(x),\\
\tilde{F}_n(\tilde{\alpha}_nx + \tilde{\beta}_n) \rightarrow G_2(x)
\end{align}
\end{proof}
MY VERSION of Khintchine's theorem:
\begin{proof}
Consider an interval $I\subset\mathbb{R}$ such that $\forall x\in I: 0 < G_{1,2}(x) < 1$.
Assume that
\begin{align}
F_n(a_nx + b_n) &\rightarrow G_1(x),\\
F_n(\alpha_nx + \beta_n) &\rightarrow G_2(x).
\end{align}
Define $\tilde{F}_n(x) := F_n(a_nx + b_n)$, therefore
\begin{align}
\tilde{F}_n(x) &\rightarrow G_1(x),\\
\tilde{F}_n(\tilde{a}_n x + \tilde{b}_n) &= F_n\big( \underbrace{a_n \tilde{a}_n}_{:=\ \alpha_n} x + \underbrace{a_n \tilde{b}_n + b_n}_{:=\ \beta_n} \big) \rightarrow G_2(x)
\end{align}
where we may identify $\tilde{a}_n = \frac{\alpha_n}{a_n} \rightarrow a$ and $\tilde{b}_n = \frac{\beta_n - b_n}{a_n} \rightarrow b$ (NOTE: sequences are bounded, hence there exist convergent subsequences. Relabel and continue)\\
But then since $\tilde{F}_n \rightarrow G_1$, we find
\begin{align}
G_1(ax + b) &= G_2(x).
\end{align}
But now for $k\in\mathbb{N}$ consider the subsequences $n = mk$ where $m\in\mathbb{N}$ (essentially we consider the subsequences $n=1,2,3,\dots$ for $k=1$, or $n=3,6,9,\dots$ for $k=3$ etc.).
Then we find 
\begin{align}
F_{mk}(a_{mk}x + b_{mk}) &\overset{m\,\rightarrow\,\infty}{\longrightarrow} G(x) =: G_1(x).
\end{align}
Since $k$ is fixed, we may define sequences $\alpha_m := a_{mk}$ and $\beta_m := b_{mk}$.
Taking the $k$-th root of the previous expression yields
\begin{align}
F_{mk}(\alpha_mx + \beta_m)^{1/k} &\overset{m\,\rightarrow\,\infty}{\longrightarrow} G(x)^{1/k}.
\end{align}
In our special case the sequences of functions are $F_n=F^n$, and thus the above simplifies to
\begin{align}
F(\alpha_mx + \beta_m)^m &\overset{m\,\rightarrow\,\infty}{\longrightarrow} G(x)^{1/k} =: G_2(x).
\end{align}
However, we may interpret this as $F_m(\alpha_mx + \beta_m) \rightarrow G_2(x)$ and thus the previous theorem implies
\begin{align}
G_2(x) &= G_1(A_kx + B_k) && \Rightarrow & G(x) &= G(A_kx + B_k)^k
\end{align}
for all $k\in\mathbb{N}$, where $A_k,B_k$ are the result of the respective limits from Khintchine's theorem.
This is also referred to as the \emph{max-stable} property and completes the proof that all extreme value distributions must be max-stable.
\end{proof}
%A non-degenerate c.d.f. $G$ is max-stable if and only if there is a sequence $F_n$ and constants $a_n>0,b_n$ such that
%\begin{align}
%F_n \kla\left( \frac{x}{a_{nk}} + b_{nk} \right) \rightarrow G(x)^{1/k}\quad (n\rightarrow\infty, k\in\mathbb{N}).
%\end{align}
%Also $D(G)$ is non-empty if and only if $G$ is max-stable and then $G\in D(G)$.
%\begin{proof}
%If $G$ is non-degenerate, so is $G^{1/k}$ for each $k$
%\end{proof}
\subsection{Max-stable property}
Furthermore, if $G$ is max-stable, there exist real functions $a(s)>0,b(s)$ defined for $s>0$ such that
\begin{align}
G(a(s)x + b(s))^s = G(x)\quad \forall x,s>0 \label{eqn:khintchine theorem corollary}
\end{align}
\begin{proof}
Since $G$ is max-stable, there exist $a_n>0, b_n$ such that 
\begin{align}
G(a_nx + b_n)^n &= G(x)
\end{align}
and hence
\begin{align}
G(a_{\lfloor ns \rfloor} x + b_{\lfloor ns \rfloor})^{\lfloor ns \rfloor} &= G(x).
\end{align}
But then taking the $\frac{1}{s}$-th power of this equation yields
\begin{align}
G(a_{\lfloor ns \rfloor}x + b_{\lfloor ns \rfloor})^{\frac{1}{s} \lfloor ns \rfloor} &= G(x)^{1/s},\\
\Rightarrow\quad G(a_{\lfloor ns \rfloor}x + b_{\lfloor ns \rfloor})^n &\rightarrow G(x)^{1/s}.
\end{align}
Khintchine's theorem applies and $\alpha_n := a_{\lfloor ns \rfloor}, \beta_n := b_{\lfloor ns \rfloor}$ leads to $G(a(s)x + b(s)) = G(x)^{1/s}$ for some $a(s) > 0, b(s)$.
\end{proof}
\subsection{Max-Stable Distributions and Extremal Type Theorem}
We say that a non-degenerate d.f. $G$ is \emph{max-stable} if $G(a_n x + b_n)^n = G(x)$ for $n\in\mathbb{N}$.
Also, $G_1$ and $G_2$ are considered to be of the same \emph{type} if $G_2(x) = G_1(ax+b)$ for some $a>0,b$.
Then we may rephrase the previous definition as ``$G$ is max-stable if $G^n$ and $G$ have the same type for $n\in\mathbb{N}$''.\\
To prove the Extremal Type Theorem we start by taking the logarithm of \autoref{eqn:khintchine theorem corollary} twice giving
\begin{align}
s\log G(a(s)x + b(s)) &= \log G(x),\\
\Rightarrow\ \ \ \log \kla\left( -\log G(a(s)x + b(s)) \right) + \log s &= \log (-\log G(x)).
\end{align}
Since $G(x)$ is increasing, so is $\log G(x)$ and therefore $\psi(x) := -\log(-\log G(x))$ is also increasing and has an inverse denoted by $U(y)$.
In this form the previous equations reads
\begin{align}
\psi(a(s)x + b(s)) - \log s &= \psi(x),
\end{align}
where we may identify $c=\log s$ from the Lemma on inverse functions.
This gives us the relation
\begin{align}
\frac{U(y + \log s) - b(s)}{a(s)} &= U(y),\\
\Rightarrow\ \ \ \frac{U(y + \log s) - U(\log s)}{a(s)} &= U(y) - U(0),\\
\Rightarrow\ \ \ \tilde{U}(y+z) - \tilde{U}(z) &= \tilde{U}(y) \tilde{a}(z) \label{eqn:extremal type theorem u tilde}
\end{align}
where we defined $z := \log s$, $\tilde{a}(z) := a(e^z)$ and $\tilde{U}(y) := U(y) - U(0)$.
Since this has to hold for arbitrary $y,z$, we may interchange them.
Then from \autoref{eqn:extremal type theorem u tilde} we find
\begin{align}
\tilde{U}(y+z) &= \tilde{U}(z) + \tilde{U}(y) \tilde{a}(z) = \tilde{U}(y) + \tilde{U}(z) \tilde{a}(y),\\
\Rightarrow\ \ \  \tilde{U}(z) \kla\left( 1 - \tilde{a}(y) \right) &= \tilde{U}(y) \kla\left( 1 - \tilde{a}(z) \right). \label{eqn:extremal type theorem u tilde v2}
\end{align}
There are now two possible scenarios (a) and (b).
In the first case (a) we assume $\tilde{a}(z) \equiv 1$, which yields $\tilde{U}(y+z) = \tilde{U}(y) + \tilde{U}(z)$ and therefore $\tilde{U}(z) \propto z$.
Then $\psi^{-1}(y) = U(y) = \rho y + U(0)$ and we find
\begin{align}
x &= \psi^{-1}(\psi(x)) = \rho \psi(x) + U(0),\\
\Rightarrow\ \ \ \psi(x) &= \frac{x - U(0)}{\rho}.
\end{align}
This leads to the d.f. of Type I
\begin{align}
G(x) &= \e^{-\e^{-\psi(x)}} = \e^{-\e^{-\frac{x - U(0)}{\rho}}}.
\end{align}
The second case (b) is that $\tilde{a}(z) \neq 1$ for at least some $z$.
Thus we may divide through by $1 - \tilde{a}(z)$ in \autoref{eqn:extremal type theorem u tilde v2} and find
\begin{align}
\tilde{U}(y) &= \frac{1 - \tilde{a}(y)}{1 - \tilde{a}(z)}\tilde{U}(z) =: c(z) \kla\left( 1 - \tilde{a}(y) \right).
\end{align}
Inserting this back into \autoref{eqn:extremal type theorem u tilde} yields
\begin{align}
c(z) \kla\left( 1 - \tilde{a}(x+y) \right) - c(z) \kla\left( 1 - \tilde{a}(x) \right) &= c(z) \kla\left( 1 - \tilde{a}(y) \right) \tilde{a}(x),\\
\Rightarrow\ \ \ \tilde{a}(x) - \tilde{a}(x+y) &= \tilde{a}(x) - \tilde{a}(x) \tilde{a}(y),\\
\Rightarrow\ \ \ \tilde{a}(x+y) &= \tilde{a}(x) \tilde{a}(y).
\end{align}
But since $\tilde{a}$ is monotone (why?) we find $\tilde{a}(x) = \e^{\rho x}$ and therefore
\begin{align}
U(y) &= U(0) + \tilde{U}(y) = U(0) + c \kla\left( 1 - \e^{\rho y} \right) \equiv B + A\e^{\rho y}.
\end{align}
Since $G(x)$ is increasing, so are $\log G, \psi = -\log (-\log G)$ and $U = \psi^{-1}$.
This restricts the signs of $A,\rho$ to only two out of four cases, since $\rho > 0$ implies $A>0$ and $\rho < 0$ implies $A<0$.
Solving for $G$ finally yields
\begin{align}
G(x) &= \e^{-\e^{-\psi(x)}} = \e^{-\e^{-\frac{1}{\rho} \log \left( \frac{x - B}{A} \right)}} = \e^{-\kla\left( \frac{x - B}{A} \right)^\alpha} = \e^{-\kla\left( -\frac{\alpha}{|\alpha|} \frac{x - B}{|A|} \right)^\alpha}
\end{align}
where $\alpha := -\frac{1}{\rho}$ must have the opposite sign of $A$.
The final expression allows us to fix $A>0$ and decomposes into the two missing extremal types for $\alpha > 0$ (Weibull) and $\alpha < 0$ (Fr\'echet).


%\begin{thebibliography}{9}
%\bibitem{example} description, \url{https://www.exmapleurl}, day.\,month~2021.
%\end{thebibliography}

\end{document}